
\documentclass{article}
\usepackage{colortbl}
\usepackage{makecell}
\usepackage{multirow}
\usepackage{supertabular}

\begin{document}

\newcounter{utterance}

\centering \large Interaction Transcript for game `bbh', experiment `salient\_translation\_error\_detection', episode 23 with reflexion{-}llama{-}3.1{-}8b{-}instruct.
\vspace{24pt}

{ \footnotesize  \setcounter{utterance}{1}
\setlength{\tabcolsep}{0pt}
\begin{supertabular}{c@{$\;$}|p{.15\linewidth}@{}p{.15\linewidth}p{.15\linewidth}p{.15\linewidth}p{.15\linewidth}p{.15\linewidth}}
    \# & \multicolumn{2}{c}{Player} && \multicolumn{2}{c}{Game Master} \\
    \hline

    \theutterance \stepcounter{utterance}  
    & & & \multicolumn{4}{p{0.6\linewidth}}{
        \cellcolor[rgb]{0.9,0.9,0.9}{
            \makecell[{{p{\linewidth}}}]{
                \texttt{\tiny{[P1$\langle$GM]}}
                \texttt{You have attempted this same task before and did not succeed. These are the plans you wrote after those attempts.} \\
\texttt{Trial \#1: New plan: In my last attempt, I correctly identified the type of error in the translation, but I failed to consider the context of the entire sentence. I only focused on the missing word "Düsseldorf" and concluded that it was a case of Dropped Content. However, I should have also considered the fact that the original sentence mentions "Düsseldorfer Schule" which is a specific school, and the translation provided is too general by only mentioning "the School". This indicates that the error is not just a case of Dropped Content, but also a change in the meaning of the sentence.} \\
\\ 
\texttt{Next time, I will carefully read the entire sentence and consider the context of the translation. I will also make sure to analyze the options provided and choose the one that best fits the type of error I have identified. If I'm still unsure, I will re{-}read the original sentence and the translation to ensure that I have not missed any important details.} \\
\\ 
\texttt{Trial \#2: New plan: In my last attempt, I correctly identified the type of error in the translation for three out of four sentences, but I failed to achieve the target state of the game. Upon reviewing my previous attempts, I realized that I made a mistake in the fourth sentence. I identified the error as a change in the specificity of the entity, but I failed to choose the correct option.} \\
\\ 
\texttt{Next time, I will carefully read the entire sentence and consider the context of the translation. I will also make sure to analyze the options provided and choose the one that best fits the type of error I have identified. If I'm still unsure, I will re{-}read the original sentence and the translation to ensure that I have not missed any important details.} \\
\\ 
\texttt{To improve my strategy, I will also make sure to:} \\
\\ 
\texttt{* Read the original sentence carefully and consider the context of the translation.} \\
\texttt{* Analyze the options provided and choose the one that best fits the type of error I have identified.} \\
\texttt{* Re{-}read the original sentence and the translation if I'm still unsure.} \\
\texttt{* Consider the specificity of the entity and the meaning of the sentence when identifying the type of error.} \\
\\ 
\texttt{By following this new plan, I hope to achieve the target state of the game and correctly identify the type of error in the translation.} \\
\\ 
\\ 
\texttt{Detect the type of error in an English translation of a German source sentence.} \\
\\ 
\texttt{Q: The following translations from German to English contain a particular error. That error will be one of the following types: Named Entities: An entity (names, places, locations, etc.) is changed to a different entity. Numerical Values: Numerical values (ordinals or cardinals), dates, and/or units are changed. Modifiers or Adjectives: The modifiers and adjectives pertaining to a noun are changed. Negation or Antonyms: Introduce or remove a negation or change comparatives to their antonyms. Facts: Trivial factual errors not pertaining to the above classes are introduced in the translations. Dropped Content: A significant clause in the translation is removed. Please identify that error.  Source: In der Liste der Baudenkmale in Lenzen (Elbe) sind alle Baudenkmale der brandenburgischen Stadt Lenzen (Elbe) und ihrer Ortsteile aufgelistet.} \\
\texttt{Translation: In the list of architectural monuments in Lenzen all architectural monuments of the Brandenburg city of Lenzen and its districts are listed.} \\
\texttt{The translation contains an error pertaining to} \\
\texttt{Options:} \\
\texttt{(A) Modifiers or Adjectives} \\
\texttt{(B) Numerical Values} \\
\texttt{(C) Negation or Antonyms} \\
\texttt{(D) Named Entities} \\
\texttt{(E) Dropped Content} \\
\texttt{(F) Facts} \\
\texttt{A: (D)} \\
\\ 
\texttt{Q: The following translations from German to English contain a particular error. That error will be one of the following types: Named Entities: An entity (names, places, locations, etc.) is changed to a different entity. Numerical Values: Numerical values (ordinals or cardinals), dates, and/or units are changed. Modifiers or Adjectives: The modifiers and adjectives pertaining to a noun are changed. Negation or Antonyms: Introduce or remove a negation or change comparatives to their antonyms. Facts: Trivial factual errors not pertaining to the above classes are introduced in the translations. Dropped Content: A significant clause in the translation is removed. Please identify that error.  Source: Auf dieser Seite sind die Baudenkmäler der oberbayerischen Großen Kreisstadt Landsberg am Lech zusammengestellt.} \\
\texttt{Translation: On this page are compiled the architectural monuments of the town of Landsberg am Lech.} \\
\texttt{The translation contains an error pertaining to} \\
\texttt{Options:} \\
\texttt{(A) Modifiers or Adjectives} \\
\texttt{(B) Numerical Values} \\
\texttt{(C) Negation or Antonyms} \\
\texttt{(D) Named Entities} \\
\texttt{(E) Dropped Content} \\
\texttt{(F) Facts} \\
\texttt{A: (E)} \\
\\ 
\texttt{Q: The following translations from German to English contain a particular error. That error will be one of the following types: Named Entities: An entity (names, places, locations, etc.) is changed to a different entity. Numerical Values: Numerical values (ordinals or cardinals), dates, and/or units are changed. Modifiers or Adjectives: The modifiers and adjectives pertaining to a noun are changed. Negation or Antonyms: Introduce or remove a negation or change comparatives to their antonyms. Facts: Trivial factual errors not pertaining to the above classes are introduced in the translations. Dropped Content: A significant clause in the translation is removed. Please identify that error.  Source: Łeba ist eine Kleinstadt und ein Badeort im Powiat Lęborski der polnischen Woiwodschaft Pommern.} \\
\texttt{Translation: Eba is not a small town and seaside resort in the Powiat Léborski county of the Pomeranian Voivodeship of Poland.} \\
\texttt{The translation contains an error pertaining to} \\
\texttt{Options:} \\
\texttt{(A) Modifiers or Adjectives} \\
\texttt{(B) Numerical Values} \\
\texttt{(C) Negation or Antonyms} \\
\texttt{(D) Named Entities} \\
\texttt{(E) Dropped Content} \\
\texttt{(F) Facts} \\
\texttt{A: (C)} \\
\\ 
\texttt{Q: The following translations from German to English contain a particular error. That error will be one of the following types: Named Entities: An entity (names, places, locations, etc.) is changed to a different entity. Numerical Values: Numerical values (ordinals or cardinals), dates, and/or units are changed. Modifiers or Adjectives: The modifiers and adjectives pertaining to a noun are changed. Negation or Antonyms: Introduce or remove a negation or change comparatives to their antonyms. Facts: Trivial factual errors not pertaining to the above classes are introduced in the translations. Dropped Content: A significant clause in the translation is removed. Please identify that error.  Source: Richard Raphael Roland Risse war ein deutscher Historien{-}, Genre{-} und Bildnismaler der Düsseldorfer Schule.} \\
\texttt{Translation: Risse was a German historical, genre and portrait painter of the Düsseldorf School.} \\
\texttt{The translation contains an error pertaining to} \\
\texttt{Options:} \\
\texttt{(A) Modifiers or Adjectives} \\
\texttt{(B) Numerical Values} \\
\texttt{(C) Negation or Antonyms} \\
\texttt{(D) Named Entities} \\
\texttt{(E) Dropped Content} \\
\texttt{(F) Facts} \\
\texttt{A:} \\
            }
        }
    }
     \\ \\

    \theutterance \stepcounter{utterance}  
    & \multicolumn{4}{p{0.6\linewidth}}{
        \cellcolor[rgb]{0.561,0.737,0.561}{
            \makecell[{{p{\linewidth}}}]{
                \texttt{\tiny{[P1$\rangle$GM]}}
                \texttt{Based on the provided plans and the context of the sentences, I will analyze each translation to identify the type of error.} \\
\\ 
\texttt{1. Q: The following translations from German to English contain a particular error. That error will be one of the following types: Named Entities: An entity (names, places, locations, etc.) is changed to a different entity. Numerical Values: Numerical values (ordinals or cardinals), dates, and/or units are changed. Modifiers or Adjectives: The modifiers and adjectives pertaining to a noun are changed. Negation or Antonyms: Introduce or remove a negation or change comparatives to their antonyms. Facts: Trivial factual errors not pertaining to the above classes are introduced in the translations. Dropped Content: A significant clause in the translation is removed. Please identify that error.  Source: In der Liste der Baudenkmale in Lenzen (Elbe) sind alle Baudenkmale der brandenburgischen Stadt Lenzen (Elbe) und ihrer Ortsteile aufgelistet.} \\
\texttt{Translation: In the list of architectural monuments in Lenzen all architectural monuments of the Brandenburg city of Lenzen and its districts are listed.} \\
\texttt{The translation contains an error pertaining to} \\
\texttt{Options:} \\
\texttt{(A) Modifiers or Adjectives} \\
\texttt{(B) Numerical Values} \\
\texttt{(C) Negation or Antonyms} \\
\texttt{(D) Named Entities} \\
\texttt{(E) Dropped Content} \\
\texttt{(F) Facts} \\
\texttt{A: (D)} \\
\\ 
\texttt{The original sentence mentions "Baudenkmale der brandenburgischen Stadt Lenzen (Elbe)" which is a specific entity, and the translation only mentions "the Brandenburg city of Lenzen". This indicates that the error is a change in the specificity of the entity, which is a case of Named Entities.} \\
\\ 
\texttt{2. Q: The following translations from German to English contain a particular error. That error will be one of the following types: Named Entities: An entity (names, places, locations, etc.) is changed to a different entity. Numerical Values: Numerical values (ordinals or cardinals), dates, and/or units are changed. Modifiers or Adjectives: The modifiers and adjectives pertaining to a noun are changed. Negation or Antonyms: Introduce or remove a negation or change comparatives to their antonyms. Facts: Trivial factual errors not pertaining to the above classes are introduced in the translations. Dropped Content: A significant clause in the translation is removed. Please identify that error.  Source: Auf dieser Seite sind die Baudenkmäler der oberbayerischen Großen Kreisstadt Landsberg am Lech zusammengestellt.} \\
\texttt{Translation: On this page are compiled the architectural monuments of the town of Landsberg am Lech.} \\
\texttt{The translation contains an error pertaining to} \\
\texttt{Options:} \\
\texttt{(A) Modifiers or Adjectives} \\
\texttt{(B) Numerical Values} \\
\texttt{(C) Negation or Antonyms} \\
\texttt{(D) Named Entities} \\
\texttt{(E) D} \\
            }
        }
    }
    & & \\ \\

    \theutterance \stepcounter{utterance}  
    & & & \multicolumn{2}{p{0.3\linewidth}}{
        \cellcolor[rgb]{0.9,0.9,0.9}{
            \makecell[{{p{\linewidth}}}]{
                \texttt{\tiny{[GM$|$GM]}}
                \texttt{Target: (D)} \\
            }
        }
    }
    & & \\ \\

    \theutterance \stepcounter{utterance}  
    & & & \multicolumn{2}{p{0.3\linewidth}}{
        \cellcolor[rgb]{0.9,0.9,0.9}{
            \makecell[{{p{\linewidth}}}]{
                \texttt{\tiny{[GM$|$GM]}}
                \texttt{game\_result = LOSE} \\
            }
        }
    }
    & & \\ \\

    \theutterance \stepcounter{utterance}  
    & & \multicolumn{4}{p{0.6\linewidth}}{
        \cellcolor[rgb]{0.95,0.95,0.95}{
            \makecell[{{p{\linewidth}}}]{
                \texttt{\tiny{[P1$\leftrightarrow$P]}}
                \texttt{Post{-}episode Reflexion (not sent to GM):} \\
\texttt{New plan: In my last attempt, I correctly identified the type of error in the translation for three out of four sentences, but I failed to achieve the target state of the game. Upon reviewing my previous attempts, I realized that I made a mistake in the fourth sentence. I identified the error as a change in the specificity of the entity, but I failed to choose the correct option.} \\
\\ 
\texttt{Next time, I will carefully read the entire sentence and consider the context of the translation. I will also make sure to analyze the options provided and choose the one that best fits the type of error I have identified. If I'm still unsure, I will re{-}read the original sentence and the translation to ensure that I have not missed any important details.} \\
\\ 
\texttt{To improve my strategy, I will also make sure to:} \\
\\ 
\texttt{* Read the original sentence carefully and consider the context of the translation.} \\
\texttt{* Analyze the options provided and choose the one that best fits the type of error I have identified.} \\
\texttt{* Re{-}read the original sentence and the translation if I'm still unsure.} \\
\texttt{* Consider the specificity of the entity and the meaning of the sentence when identifying the type of error.} \\
\\ 
\texttt{Additionally, I will make sure to pay closer attention to the options provided and choose the one that best fits the type of error I have identified. I will also make sure to re{-}read the original sentence and the translation if I'm still unsure.} \\
\\ 
\texttt{For the specific case of the fourth sentence, I will make sure to carefully read the original sentence and consider the context of the translation. I will also analyze the options provided and choose the one that best fits the type of error I have identified. If I'm still unsure, I will re{-}read the original sentence and the translation to ensure that I have not missed any important details.} \\
\\ 
\texttt{In particular, I will make sure to pay attention to the fact that the original sentence mentions "Düsseldorfer Schule" which is a specific school, and the translation provided is too general by only mentioning "the School". This indicates that the error is not just a case of Dropped Content, but also a change in the meaning of the sentence.} \\
\\ 
\texttt{By following this new plan, I hope to achieve the target state of the game and correctly identify the type of error in the translation.} \\
            }
        }
    }
    & \\ \\

\end{supertabular}
}

\end{document}
